% !TEX root = ../glossario.tex
% Lettera P del Glossario (vedi capitoli/a.tex per le istruzioni).
\phantomsection
\letteraglossario{P}

\termine{PoC}{\emph{Proof of Concept}: prototipo esplorativo che dimostra la
  fattibilità delle principali scelte tecnologiche; nel progetto
  didattico è presentato alla Requirements and Technology Baseline
  (RTB).}

\termine{Prompt}{Testo o istruzione formulata in linguaggio naturale che viene fornita in input a un modello linguistico (LLM) per ottenere una risposta o fargli eseguire un'operazione.}

\termine{Proponente}{Organizzazione che propone un progetto software, fornendo i requisiti e le specifiche di alto livello, e che ne verifica la realizzazione.}

\termine{Pull request}{Richiesta formale (tipica delle piattaforme come GitHub) per integrare le modifiche effettuate su un ramo (branch) di sviluppo all'interno del ramo principale, favorendo la revisione del codice (code review) prima dell'unione.}

\termine{Push}{Comando del sistema di controllo di versione Git utilizzato per inviare le modifiche confermate (commit) dal repository locale a quello remoto.}

\termine{PWA}{\emph{Progressive Web App}: applicazione realizzata con tecnologie
  web, accessibile da browser e installabile sui dispositivi
  compatibili come un'app nativa.}
